\chapter{LITERATURE SURVEY}\label{ch:literature}

\section{Literature Review}\label{sec:literature-review}
\textbf{X. Zhuang et al.,} presents an innovative approach to military target detection by combining EfficientDet with a Generative Adversarial Network to improve training data quality and detection accuracy. Initially, EfficientDet is applied to identify targets within the original dataset. Subsequently, the GAN generates synthetic military target images, which are incorporated into the training dataset to enhance its diversity and robustness. The expanded dataset is then used to retrain the EfficientDet model, leading to significantly improved detection performance, particularly in situations where limited training data is available. Experimental evaluation on a custom military target dataset demonstrated the effectiveness of the proposed method. The model achieved a mean Average Precision (mAP) of 92.8\%, surpassing the performance of conventional approaches such as YOLOv3 and YOLOv5. Moreover, the system maintained real-time processing capability at 15 FPS and achieved over 90\% detection accuracy under challenging conditions, including dim lighting and occlusions. \textbf{\cite{base_paper}} \\

\textbf{C. Konar et al.,} presents a portable and edge-optimized surveillance system developed to detect modern security threats such as concealed weapons and drones. The proposed framework is designed for deployment in resource-constrained environments, including remote border regions and civilian infrastructure, while maintaining military-grade object detection capabilities. The system is built upon the Florence v2 Vision-Language Transformer and incorporates Parameter Efficient Fine-Tuning (PEFT) along with Low-Rank Adaptation to improve adaptation efficiency with reduced computational overhead. The model was trained using a dataset of 12,000 images, including GAN-generated synthetic samples to enhance detection robustness. For deployment, the framework utilized a Raspberry Pi 5 integrated with an ESP32-S2 module, enabling wireless alert transmission through the ESP-NOW communication protocol. Experimental results demonstrated the effectiveness of the proposed system, achieving an accuracy of 88.01\% and outperforming YOLOv8 and Faster R-CNN. Additionally, the model attained a mAP$_{50}$ score of 64.83 while maintaining high efficiency on edge hardware with low CPU and memory utilization. \textbf{\cite{supplementary_paper_1}} \\

\textbf{A. A. S. Matta et al.,} presents an optimized spatiotemporal deep learning framework for predictive behavioral threat detection in surveillance footage. The proposed system addresses challenges in anomaly detection, including complex human behavior, temporal variability, and limited annotated data, by integrating a Convolutional Neural Network with Long Short-Term Memory networks. The CNN component extracts spatial features from video frames, while the LSTM captures temporal dependencies across sequential frames, enabling effective recognition of suspicious activities. The framework also incorporates hyperparameter optimization, regularization techniques, and cross-validation strategies to improve model stability and generalization. Experimental evaluation was conducted using the DCSASS surveillance dataset containing multiple anomaly categories such as robbery, assault, and vandalism. Results demonstrated that the optimized CNN-LSTM model achieved an impressive accuracy of 98.1\%, with high precision, recall, and F1-scores across 3-fold, 5-fold, and 10-fold validation settings. Furthermore, the proposed framework outperformed traditional machine learning models and recent deep learning approaches, proving highly effective for real-time intelligent surveillance applications. \textbf{\cite{supplementary_paper_2}} \\

\textbf{N. Li et al.,} proposes PDGV-DETR, an advanced framework for precise weapon and personnel detection in static security surveillance images. The study addresses critical challenges such as object occlusion and complex backgrounds, which often reduce the effectiveness of conventional detection models in public safety applications, including school and subway surveillance systems. The proposed framework enhances detection performance through three major components. First, a dynamic hierarchical channel interaction convolution module improves the recognition of partially occluded objects. Second, an improved bidirectional hybrid feature pyramid network enables effective cross-scale feature fusion for detecting objects of varying sizes. Finally, a global semantic weaving network strengthens the model's ability to distinguish targets from cluttered and complex backgrounds in static scenes. Experimental evaluation demonstrated the effectiveness of PDGV-DETR across multiple datasets. The framework achieved a peak mAP$_{50}$ of 85.9\% on conflict-scene datasets and 93\% on weapon-specific datasets. Furthermore, the model significantly outperformed RT-DETR and Deformable DETR, while maintaining a balance between high localization accuracy and computational efficiency suitable for real-time surveillance applications. \textbf{\cite{supplementary_paper_3}} \\

\textbf{A. M. Garcia-Rubio et al.,} presents a collaborative AI-human framework for detecting concept drift in object detection models used in defense and surveillance applications. The proposed system addresses the degradation of object detection performance caused by evolving environments, lighting conditions, and unfamiliar objects. The framework integrates a YOLOv7 architecture with human feedback to improve reliability and adaptability in dynamic military scenarios. A custom metric called the Drift Avoidance Value was introduced to quantify concept drift by combining AI-generated confidence scores with human corrections, duplicate detection analysis, and missed object evaluations. The model was trained on 4,579 military vehicle images and evaluated using multiple datasets containing main battle tanks, humans, and unknown ground vehicles. Experimental results demonstrated that retraining the model with human-corrected data significantly improved detection performance. The Drift Avoidance Value increased after retraining, while prediction accuracy improved by 16.9\% for tanks and 12.5\% for humans. The study highlights the effectiveness of human-in-the-loop learning for maintaining robust and adaptive object detection systems in real-world defense environments. \textbf{\cite{supplementary_paper_4}} \\

\textbf{X. Xiang et al.,} introduces UniFusOD, an end-to-end framework designed for infrared--visible image fusion and object detection. The proposed system combines visible images, which provide rich texture and color information, with infrared images that offer strong edge details and reliable performance under low-light conditions. By integrating these complementary modalities, UniFusOD generates a single fused image optimized for both visual clarity and accurate object detection. The framework incorporates the FRA mechanism, which utilizes multiple convolution kernels to capture fine-grained details as well as broader semantic features. Additionally, the proposed UnityGrad Optimization strategy balances fusion and detection gradients, preventing conflicts between the two tasks and enabling more effective joint training. Experimental evaluation demonstrated that UniFusOD outperformed several state-of-the-art methods across multiple benchmark datasets. The model achieved improved mean Average Precision (mAP) scores, including gains of 0.8 on the DroneVehicle Dataset and 1.9 on the M3FD Dataset. Furthermore, the framework produced fused images with superior structural quality and information content while maintaining strong robustness in detecting small targets within complex environments. \textbf{\cite{supplementary_paper_5}} \\

\textbf{Z. Shou et al.,} presents a robust framework for cyber physical systems, including autonomous vehicles and smart surveillance applications, through the development of MSAF-Net for human action recognition. The proposed system integrates multiple visual modalities such as RGB, depth, and thermal images to enhance detection accuracy and reliability in complex environments. The framework employs a Multi-Scale Attention Fusion mechanism that extracts features at different levels of detail while emphasizing the most relevant information for accurate recognition. In addition, Cross-Level Feature Interaction enables different network layers to exchange information, preserving both fine-grained details and global contextual structures. The model also incorporates an Adaptive Fusion Strategy that combines semantic and structural guidance to generate clear and consistent fused representations. Experimental evaluation on several benchmark datasets demonstrated the effectiveness of the proposed approach. MSAF-Net achieved 91.54\% accuracy on the UCF101 Dataset and 92.14\% accuracy on the ActivityNet, outperforming existing state-of-the-art methods. Furthermore, the framework remained robust even when individual sensors failed and could be efficiently pruned for deployment on low-power hardware platforms. \textbf{\cite{supplementary_paper_6}} \\

\textbf{Z. Song et al.,} introduces EDNet, a deep learning framework designed for detecting small objects in UAV and drone imagery. Built upon YOLOv10, the framework is optimized for edge computing, enabling efficient deployment on mobile and low-power devices. EDNet is highly scalable and offers seven model variants ranging from Tiny, which is lightweight and computationally efficient, to XL, which provides high-performance detection capabilities. The architecture incorporates the C2f-FCA block, which utilizes faster context attention mechanisms to extract meaningful features while maintaining low computational complexity. In addition, the model employs the WIoUv3 Loss Function to prioritize high-quality samples and reduce the influence of noisy data during training. For optimized deployment, the framework supports technologies such as CoreML, TensorRT, and OpenVINO, ensuring high-speed performance across multiple hardware platforms. Experimental results demonstrated that EDNet outperformed YOLOv8 and YOLOv10 in both accuracy and efficiency, achieving up to 5.6\% higher mAP$_{50}$ with fewer parameters while maintaining real-time performance across devices such as the iPhone 12 and Raspberry Pi 5. \textbf{\cite{supplementary_paper_7}} \\

\textbf{S. D N et al.,} presents a state-of-the-art object detection system designed for analyzing camera and drone footage in military environments. The proposed framework focuses on identifying battlefield threats such as vehicles, weapons, and combatants in real time. The system employs a vision-language model that combines image data with natural language descriptions to improve object understanding and classification accuracy. It integrates CLIP to associate images with textual descriptions, enabling the detection of unseen objects without requiring extensive labeled datasets. Additionally, the SPPELAN pooling module is used to generate fixed feature representations while minimizing information loss. The framework also incorporates an N-gram and GLIP labeler to extract noun phrases and create pseudo-boundaries for automated training data generation. Furthermore, K-means Clustering is applied to group similar image regions, improving detection efficiency in cluttered environments. Experimental evaluation on the Microsoft COCO Dataset demonstrated superior performance compared to Grounding DINO, achieving lower complexity, higher accuracy, and significantly faster frame rates suitable for real-time deployment. \textbf{\cite{supplementary_paper_8}} \\

\textbf{V. Kaya et al.,} proposes a deep learning-based approach for the automatic detection and classification of seven different handheld weapon types. The study aims to assist security agencies by developing a proactive surveillance system capable of identifying threatening objects in real time, thereby helping to prevent criminal activities. The proposed model is based on the VGGNet architecture and consists of 25 layers with 337,671 parameters, making it both efficient and computationally lightweight. A custom dataset containing 5,214 images of various weapons, including assault rifles, pistols, and grenades, was collected from internet sources and preprocessed using grayscale conversion and background cleaning techniques. The model was implemented and trained using the Keras deep learning framework. Experimental results demonstrated that the proposed CNN model achieved an impressive classification accuracy of 98.40\%, significantly outperforming models such as VGG-16, ResNet50, and ResNet-101. Additionally, the model achieved a low loss rate of 0.052, highlighting its reliability, fast training capability, and suitability for deployment on low-cost computing systems. \textbf{\cite{supplementary_paper_9}} \\

\textbf{Ankita et al.,} introduces an efficient and lightweight deep learning framework for human activity recognition using smartphone sensor data. The study addresses the limitations of traditional pattern recognition techniques by employing automated feature extraction, which is highly beneficial for applications such as healthcare, elderly assistance, and sports science. The proposed model integrates a Convolutional Neural Network with Long Short-Term Memory units to form a CNN-LSTM architecture capable of learning both spatial and temporal features. Convolutional layers automatically extract meaningful patterns from accelerometer and gyroscope signals, while LSTM layers analyze sequential dependencies in human activities. The model was trained and evaluated using the UCI-HAR Dataset collected from Samsung Galaxy S2 smartphones, with Gaussian standardization applied during preprocessing. Experimental results demonstrated that the CNN-LSTM model achieved an impressive recognition accuracy of 97.89\%. Furthermore, the model outperformed traditional machine learning approaches by effectively recognizing activities such as walking, sitting, and standing while maintaining a lightweight structure suitable for mobile deployment. \textbf{\cite{supplementary_paper_10}}

\label{end:literature-review}
\clearpage
\section{Comparative Study Of References}\label{sec:comparison}
\renewcommand{\arraystretch}{1.35}
\setlength{\tabcolsep}{5pt}
\begin{longtable}{|>{\centering\arraybackslash}p{1.2cm}|p{6.3cm}|p{6.3cm}|}
  \caption{Comparison Table of Research Papers}\label{tab:comparison} \\
  \hline
  \textbf{Sl. No.} & \textbf{Research Paper} & \textbf{Learning Outcome} \\
  \hline
  \endfirsthead
  \hline
  \textbf{Sl. No.} & \textbf{Research Paper} & \textbf{Learning Outcome} \\
  \hline
  \endhead
  \hline
  \multicolumn{3}{|r|}{\textit{Continued on next page}} \\
  \hline
  \endfoot
  \hline
  \endlastfoot

  1 &
  \textbf{Title:} Military target detection method based on EfficientDet and Generative Adversarial Network \par
  \textbf{Authors:} X. Zhuang, D. Li, Y. Wang, and K. Li \par
  \textbf{Published:} Engineering Applications of Artificial Intelligence, vol. 132, Jun. 2024 &
  \textbf{Objective:} To detect targets in low visibility, noisy, or data-scarce combat zones \par
  \textbf{Methodology:} EfficientDet-D0 + Dual-Discriminator GAN \par
  \textbf{Core Idea:} Uses GANs to generate synthetic military training data and BiFPN for feature fusion \\
  \hline

  2 &
  \textbf{Title:} Edge-Optimized Threat Detection with Florence v2 \par
  \textbf{Authors:} C. Konar, D. Das, D. S., K. Jc, and S. Singh \par
  \textbf{Published:} PeerJ Computer Science, vol. 12, Mar. 2026 &
  \textbf{Objective:} To create a portable edge-optimized surveillance system capable of detecting military-relevant threat classes in resource-constrained environments \par
  \textbf{Methodology:} Florence v2 Vision-Language Transformer \par
  \textbf{Core Idea:} Leveraging multimodal vision-language models on compact edge hardware to provide situational awareness without relying on cloud infrastructure \\
  \hline

  3 &
  \textbf{Title:} Spatiotemporal deep learning framework for predictive behavioral threat detection in surveillance footage \par
  \textbf{Authors:} A. A. S. Matta and V. P. C. S. R. Manukonda \par
  \textbf{Published:} Frontiers in Big Data, vol. 9, Feb. 2026 &
  \textbf{Objective:} To detect anomalous human behaviors in untrimmed surveillance footage using optimized spatiotemporal deep learning \par
  \textbf{Methodology:} Hybrid CNN-LSTM framework on DCSASS dataset \par
  \textbf{Core Idea:} Integrating CNN spatial features and LSTM temporal patterns to reliably identify 13 distinct types of criminal behavior \\
  \hline

  4 &
  \textbf{Title:} PDGV-DETR: Object Detection for Secure On-Site Weapon and Personnel Location \par
  \textbf{Authors:} N. Li et al. \par
  \textbf{Published:} Sensors, vol. 26, Feb. 2026 &
  \textbf{Objective:} To precisely detect and position weapons and personnel in static surveillance images while overcoming occlusion and complex background interference \par
  \textbf{Methodology:} PDGV-DETR framework \par
  \textbf{Core Idea:} Transformer-based end-to-end detection system that uses cross-scale semantic fusion to bridge the gap between high-level semantics and low-level image details \\
  \hline

  5 &
  \textbf{Title:} Detecting Concept Drift in Object Detection Models: A Collaborative AI-Human Approach for Defense Applications \par
  \textbf{Authors:} A. M. Garcia-Rubio, L. Heck, K. M. Schweitzer, R. M. Bateman, A. Alaeddini, and K. K. Castillo-Villar \par
  \textbf{Published:} IEEE Access, vol. 14, Jan. 2026 &
  \textbf{Objective:} To quantify and manage model performance deterioration in military object detection systems over time \par
  \textbf{Methodology:} YOLO confidence scores + direct human feedback \par
  \textbf{Core Idea:} Enhancing AI reliability by integrating human corrections directly into drift calculations to guide targeted and efficient model retraining \\
  \hline

  6 &
  \textbf{Title:} Infrared-Visible Image Fusion Meets Object Detection \par
  \textbf{Authors:} X. Xiang et al. \par
  \textbf{Published:} Remote Sensing, vol. 17, Nov. 2025 &
  \textbf{Objective:} To solve gradient conflicts when training fusion and detection together \par
  \textbf{Methodology:} UniFusOD \par
  \textbf{Core Idea:} Use Fine-Grained Region Attention (FRA) mechanism and a unified framework to fuse RGB and infrared images for accurate detection \\
  \hline

  7 &
  \textbf{Title:} Multi-modal action recognition via advanced image fusion techniques \par
  \textbf{Authors:} Z. Shou and D. Zhu \par
  \textbf{Published:} Frontiers in Physics, Aug. 2025 &
  \textbf{Objective:} To recognize human or combatant actions accurately even if one sensor fails \par
  \textbf{Methodology:} MSAF-Net (Multi-Scale Attention-Guided Fusion) \par
  \textbf{Core Idea:} Dynamically prioritizes sensor data and switches to backup sensors when needed \\
  \hline

  8 &
  \textbf{Title:} EDNet: Edge-Optimized Small Target Detection in UAV Imagery \par
  \textbf{Authors:} Z. Song, Y. Zhang, and A. A. R. M. A. Ebayyeh \par
  \textbf{Published:} IEEE SWC, Dec. 2024 &
  \textbf{Objective:} To enable high-speed detection of tiny objects on low-power drone hardware \par
  \textbf{Methodology:} YOLOv10 + Faster Context Attention (C2f-FCA) \par
  \textbf{Core Idea:} Uses an XSmall detection layer for small object detection with hardware acceleration optimization \\
  \hline

  9 &
  \textbf{Title:} Real-Time Versatile Object Detection With High Accuracy And Reduced Information Loss During Wartime \par
  \textbf{Authors:} S. D. N. and D. M. \par
  \textbf{Published:} IEEE Conference, Sep. 2024 &
  \textbf{Objective:} To detect unseen or new battlefield threats without manual re-labeling \par
  \textbf{Methodology:} CLIP Encoder + Modified C2f-Darknet \par
  \textbf{Core Idea:} Uses pseudo-boundary labeling for zero-shot detection and combines vision and language \\
  \hline

  10 &
  \textbf{Title:} Detection and Classification of Different Weapon Types Using Deep Learning \par
  \textbf{Authors:} V. Kaya, S. Tuncer, and A. Baran \par
  \textbf{Published:} Applied Sciences, vol. 11, Aug. 2021 &
  \textbf{Objective:} To propose a model for the automatic detection and classification of distinct types of handheld weapons \par
  \textbf{Methodology:} CNN based on VGGNet architecture \par
  \textbf{Core Idea:} Utilizing reduced layer complexity and high-precision CNNs to enable weapon identification on low-cost computing systems \\
  \hline

  11 &
  \textbf{Title:} An Efficient and Lightweight Deep Learning Model for Human Activity Recognition Using Smartphones \par
  \textbf{Authors:} Ankita, S. Rani, H. Babbar, S. Coleman, A. Singh, and H. M. Aljahdali \par
  \textbf{Published:} Sensors, vol. 21, Jun. 2021 &
  \textbf{Objective:} To develop an efficient and lightweight model for identifying human activities using smartphone-based sensor data \par
  \textbf{Methodology:} CNN + Long Short-Term Memory (LSTM) \par
  \textbf{Core Idea:} Integrating spatial and temporal deep learning architectures to automate feature extraction from raw sensor data \\
  \hline
\end{longtable}
\label{end:comparison}
\label{end:literature}
